\providecommand{\toplevelprefix}{../..}  %
\documentclass[../../book-main.tex]{subfiles}

\begin{document}

\chapter{Open Problems and Directions about Intelligence}
\label{ch:future}

\begin{quote}
    ``\textit{The study is to proceed on the basis of the conjecture
        that every aspect of learning or any other feature of
        intelligence can in principle be so precisely described that a
        machine can be made to simulate it. An attempt will be made to
        find how to make machines use language, form abstractions and
        concepts, solve kinds of problem now reserved for humans, and
    improve themselves.}''

    $~$\hfill -- Proposal for the Dartmouth AI program, 1956
\end{quote}
\vspace{5mm}

This manuscript systematically introduces mathematical principles and
computational mechanisms for how memory or empirical knowledge can be developed from observed high-dimensional data. 
The ability to seek parsimony in a seemingly random world is a fundamental characteristic of any
intelligence, natural or artificial. We believe the principles and
mechanisms presented in this book are unifying and universal, applicable to both animals and machines.

We hope this book helps readers fully clarify the mystery surrounding
modern practices of artificial deep neural networks by developing a
rigorous understanding of their functions and roles in learning (representations of) 
low-dimensional distributions from high-dimensional data. With this
understanding, the capabilities and limitations of existing AI models
and systems become clear:
\begin{enumerate}
    \item Existing open-loop models and systems fall short of being complete memory systems capable 
    of self-learning and self-improving.
    \item Existing realizations of representation learning remain primitive and brute-force, 
    far from optimal in terms of network architectures or optimization strategies.
    \item Existing AI models only learn data distributions and
    conduct inductive Bayesian inference, which differs from high-level human intelligence.
\end{enumerate}

One goal of this book is to help readers establish an objective and
systematic understanding of current machine intelligence technologies
and to recognize what open problems and challenges remain for further
advancement of machine intelligence. In the last chapter, we provide
some of our views and projections for the future.

\section{Towards Autonomous Intelligence: Close the Loop?}
From the practice of machine intelligence in the past decade, it has
become clear that, given sufficient data and computational resources,
one can build a large enough model and pre-train it to learn the
\textit{prior} distribution of all the data, say \(p(\vx)\).
Theoretically, such a large model can memorize almost all existing
knowledge about the world that is encoded in data such as 2D images,
3D shapes, dynamical motions, and natural languages. As we discussed
at the beginning of the book, such a large model plays a role similar
to DNA, which life uses to record and pass on knowledge about the world.

The model and distribution learned in this way can then be used to
create new data samples drawn from the same distribution. One can
also use the model to conduct inference (e.g., completion,
estimation, and prediction) with the memorized knowledge under
various conditions, say by sampling the \textit{posterior}
distribution \(p(\vx\mid \vy)\) under a new observation \(\vy\).
Strictly speaking, such inference is inductive or statistical.

Any pre-trained model, however large, cannot guarantee that the
distribution it has learned so far is entirely correct or complete.
If our samples \(\hat{\vx}_{t}\) from the current \textit{prior}
\(p_{t}(\vx)\) or estimates \(\hat{\vx}_{t}(\vy)\) based on the
\textit{posterior} \(p_{t}(\vx\mid \vy)\) are inconsistent with the
truth \(\vx\), we would like to correct the learned distributions:
\begin{equation}
    p_{t}(\vx) \rightarrow p_{t+1}(\vx) \quad \text{or}\quad
    p_{t}(\vx\mid \vy) \rightarrow p_{t+1}(\vx\mid \vy),
\end{equation}
based on the error \(\ve_{t} = \vx_{t} - \hat{\vx}_{t}\). This is
known as error correction based on error feedback, a ubiquitous
mechanism in nature for continuous learning. However, any open-ended
model itself lacks the mechanism to revise or improve the learned
distribution when it is incorrect or incomplete. Improving current AI
models still depends largely on human involvement: supervision or
reinforcement through experimentation, evaluation, and selection. We
may call this process ``artificial selection'' of large models, as
opposed to the natural selection for the evolution of life.

As we studied earlier in this book (\Cref{ch:consistent-representations} in
particular), closed-loop systems align an internal representation
with sensed observations of the external world. They can continuously
improve the internally learned distribution and its representation to
(incrementally) achieve consistency or self-consistency. An immediate next step is to
develop truly closed-loop memory systems, illustrated in
\Cref{fig:open-to-closed}, that autonomously and continuously learn
general data distributions and improve based on error feedback.

Therefore, the transition from the currently popular end-to-end
trained open-loop models to continuously learning closed-loop systems
\begin{equation}
    \text{\textbf{open-ended} models} \; \Longrightarrow \;
    \text{\textbf{closed-loop} systems}
\end{equation}
is the key for machines to truly emulate how the animal brain learns
and applies knowledge in an open world. We believe that
\begin{quote}
    \begin{center}
        \textit{open-ended models are for a closed world, however
        large; \\ closed-loop systems are for an open world, however small.}
    \end{center}
\end{quote}
Hence, the so-called ``artificial general intelligence'' (AGI) can never be 
achieved by simply having a system to memorize all existing data and knowledge 
of the world. No knowledge, however much, is truly generalizable.\footnote{According to Sir Karl Popper, 
an influential Philosopher of Science, all scientific theory and knowledge are falsifiable!} Nevertheless, 
it is the ability to improve one's existing memory or knowledge and adapt to any new environments that is 
truly generalizable. Hence, truly genuine intelligence is itself generalizable if realized 
completely and correctly.\footnote{In our opinion, AGI should represent ``artificial genuine intelligence.''} 
This is precisely the essence of the Cybernetics program laid out by Norbert Wiener in the
1940s that we discussed at the very beginning of this book.

\begin{figure}[t]
    \centering
    \includegraphics[width=0.9\linewidth]{\toplevelprefix/chapters/chapter9/figs/open-to-closed.png}
    \caption{From an open-ended deep network to a closed-loop system.}
    \label{fig:open-to-closed}
\end{figure}

\section{Towards Natural Intelligence: Beyond Back Propagation?}
The practice of machine intelligence in recent years has led many to
believe that one must build a single large model to learn the
distribution of all data and memorize all knowledge. Even though this
is technologically possible, such a solution is likely far from
necessary or efficient. As we know from training deep networks, the
only known scalable method to train such networks at scale is through
back propagation (BP) \citep{Back-Prop}. Although BP offers a way to
correct errors via gradient signals propagated back through the whole
model, it is nevertheless rather brute-force and differs
significantly from how nature learns: BP is an option that nature
cannot afford due to its high cost and simply cannot implement due to
physical limitations.

More generally, we cannot truly understand intelligence unless we
also understand how it can be efficiently implemented. That is, one
needs to address the computational complexity of realizing mechanisms
associated with achieving the objectives of intelligence.
Historically, our understanding of (machine) intelligence has evolved
through several phases, from the incomputable Kolmogorov complexity
to Shannon's entropy, from Turing's computability to later
understanding of tractability,\footnote{We say a problem is tractable
    if it allows an algorithm whose complexity is polynomial in the size
of the problem.} and to the strong emphasis on algorithm scalability
in modern practice of high-dimensional data analysis
\citep{Wright-Ma-2022} and artificial intelligence. This evolution can
be summarized as follows:
\begin{equation}
    \text{\textbf{incomputable}} \;
    \Longrightarrow \; \text{\textbf{computable}} \;
    \Longrightarrow \; \text{\textbf{tractable}} \; \Longrightarrow \;
    \text{\textbf{scalable}}.
\end{equation}
To a large extent, the success and popularity of deep learning and
back propagation is precisely because they have offered a scalable
implementation with modern computing platforms (such as GPUs) for
processing and compressing massive data. Nevertheless, such an
implementation is still far more expensive than how nature realizes
intelligence.

There remains significant room to improve the efficiency of machine
intelligence so that it can emulate the efficiency of natural
intelligence, which should be orders of magnitude greater than
current brute-force implementations. To this end, we need to discover
new learning architectures and optimization mechanisms that enable
learning data distributions under natural physical conditions and
resource constraints, similar to those faced by intelligent beings in
nature---for example, without accessing all data at once or updating
all model parameters simultaneously (via BP).

The principled framework and approach laid out in this book can guide
us to discover such new architectures and mechanisms. These new
architectures and mechanisms should enable online continuous learning
and should be updatable through highly localized and sparse forward
or backward optimization.\footnote{For learning a distribution, a
    simple instantiation of these desiderata is in the simplest case of
    PCA, with the online PCA method introduced in
    \Cref{ch:consistent-representations}. Once the linear model for the 
    distribution is learned, the easier task of conducting online prediction 
    with noisy observations can then be easily accomplished by least-square
    type methods, such as the Wiener filter \cite{Wiener-1942,Wiener-1949}. 
    In the case of linear dynamical systems, that becomes the famous
    the Kalman filter, see \cite[Appendix B]{MaY2003} or \citep{peng2025mathematics}.}

\begin{figure}[t]
    \centering
    \includegraphics[width=0.99\linewidth]{\toplevelprefix/chapters/chapter9/figs/loops.png}
    \caption{Conjectured architecture of the brain cortex. The cortex
        is a massively parallel and distributed auto-encoding system
        consisting of a hierarchy of closed-loop auto-encoders. They
        extract information from their inputs and maximize the
        information gain of the representations they output. Such processes are done at multiple
    levels of the hierarchy with different levels of abstraction of the information.}
    \label{fig:loops}
\end{figure}
As we have learned from neuroscience, the cortex of our brain
consists of tens of thousands of cortical columns
\citep{Hawkins-2021}. All cortical columns have similar physical
structures and functions. They are highly parallel and distributed,
though sparsely interconnected. Hence, we believe that to develop a
more scalable and structured memory system, we need to consider
architectures that emulate those of the cortex. \Cref{fig:loops}
shows such a hypothesized architecture: a massively distributed and
hierarchical system consisting of many largely parallel closed-loop
auto-encoding modules.\footnote{Elements of such hypothetical
    architectures exist in the literature in various forms, such as
    predictive coding networks \citep{millidge2022predictive}, which have
an optimization strategy that utilizes efficient local updates.}
These modules learn to encode different sensory modalities or many
projections of data from each sensory modality. Our discussion in
\Cref{sec:measurement-self-consistency} of
\Cref{ch:conditional-inference} suggests that such parallel sensing
and learning of a low-dimensional distribution is theoretically
possible. Higher-level (lossy) autoencoders can then be learned based
on outputs of lower-level ones to develop more sparse and
higher-level ``abstractions'' of the representations learned by the
lower levels.

The distributed, hierarchical, and closed-loop system architecture
illustrated in \Cref{fig:loops} shares many characteristics of the
brain's cortex. Such a system architecture may open up many more
possibilities than the current single large-model architecture. It
enables exploration of much more efficient learning and optimization
mechanisms and results in a more structured modular organization of
the learned data distribution and knowledge, more closely emulating the memory
in our brain. This would allow us to bring the implementation of machine 
intelligence to the next level of evolution:
\begin{equation}
    \text{\textbf{incomputable}}
    \Longrightarrow \text{\textbf{computable}}
    \Longrightarrow \text{\textbf{tractable}} \Longrightarrow
    \text{\textbf{scalable}}
    \Longrightarrow
    \text{\textbf{natural}}.
\end{equation}

\section{Towards Scientific Intelligence: Beyond the Turing Test?}

\subsection{The Evolution of Intelligence in Nature}
As we have discussed at the beginning of this book, \Cref{ch:intro},
intelligence in nature has evolved through multiple phases and
manifested in four distinct forms:
\begin{equation}
    \text{\textbf{phylogenetic}} \;
    \Longrightarrow \; \text{\textbf{ontogenetic}} \;
    \Longrightarrow \; \text{\textbf{societal}} \;
    \Longrightarrow \; \text{\textbf{scientific intelligence}}.
\end{equation}
All forms of intelligence share the common objective (and
characteristic) of learning useful knowledge as low-dimensional
distributions of sensed high-dim data about the world.
Nevertheless, they differ in the specific code books or coding schemes adopted, the
information encoded, the computational mechanisms for learning and
improving, and the physical implementations of such mechanisms. Using
the concepts and terminology developed in this book, the four stages
of intelligence differ in the following three aspects:
\begin{enumerate}
    \item The \textit{code book} or scheme used to encode the
        intended information or knowledge.
    \item The \textit{information} or knowledge learned and encoded
        using the code book.
    \item The \textit{optimization mechanisms} used to create and
        correct the encoded information or knowledge.
\end{enumerate}
\Cref{tab:four-stage-intelligence} summarizes their main characteristics.
\begin{table}[h]
    \centering
    \resizebox{1\columnwidth}{!}
    {
        \begin{tabular}{| c | c | c | c | c |}
            \hline & \textbf{Phylogenetic} & \textbf{Ontogenetic} &
            \textbf{Societal} & \textbf{Scientific}\\
            \hline

            \textbf{Codebook}  & DNAs & Neurons/Brain & Natural
            Languages & Math Abstractions \\ [0.5ex]
            \hline

            \textbf{Information} & Genes & Memory & Empirical
            Knowledge & Scientific Facts\\ [0.5ex]
            \hline
            \textbf{Optimization} & Reinforcement & Error Feedback &
            Trial \& Error & Theorize \& Falsify \\  [0.5ex]
            \hline
        \end{tabular}
    }
    \caption{Main characteristics of the four stages of intelligence in nature.}
    \label{tab:four-stage-intelligence}
\end{table}

As we now know, humans have achieved two quantum leaps in intelligence.
\begin{enumerate}
    \item The first was the development of spoken and written
        language, which enabled humans to share and transmit learned
        knowledge across generations, much as DNA does in nature.
    \item The second was the development of mathematics and formal
        logic roughly three thousand years ago, which became the
        precise language of modern science.
\end{enumerate}
The language of mathematics freed us from summarizing knowledge from
observations only in empirical form and allowed us to formalize
knowledge as theories verifiable or falsifiable through mathematical 
deduction or experimentation. Through hypothesis formulation, logical
deduction, and experimental verification or falsification, we can now
proactively discover and develop new knowledge that is far beyond
what can be learned from the distributions of observed
data. For example, causal relationships cannot be learned
from distributions alone \citep{Pearl-2009}. 

Despite the seeming differences among these different stages of
intelligence, the evolution of intelligence shares a common
characteristic: \textit{they all use certain forms of
    feedback\footnote{Loss for error and reward for success.} to create
    and correct information learned, and the feedback and correction are
    becoming increasingly frequent and efficient through evolution---from 
    open-loop reinforcement to closed-loop self-supervision, from
collective trial and error to proactive hypothesizing and falsification.}

\subsection{From Inductive to Deductive Intelligence}
As discussed in the introduction (\Cref{ch:intro}), the 1956
``artificial intelligence'' (AI) program aimed precisely to study how
to endow machines with scientific intelligence, i.e., high-level
functions such as mathematical abstraction, causal inference, logical deduction, and
problem solving that are believed to differentiate humans from animals:
\begin{equation}
    \text{\textbf{low-level} (animal) intelligence} \; \Longrightarrow \;
    \text{\textbf{high-level} (human) intelligence.}
\end{equation}
As we have clarified repeatedly in this book, most technological
advances in machine intelligence over the past decade or so, although
carried out under the name ``AI'', are actually more closely related
to having machines imitate low-level forms of intelligence largely
shared by both animals and humans, including phylogenetic,
ontogenetic, and societal intelligence. At these levels, intelligence
creates memory or \textit{empirical} knowledge (from empirical data distributions) and conducts inference
that is mostly \textit{inductive} (of Bayesian-like) in nature. Principles and methods introduced in this book aim to reveal scientific objectives and
computational mechanisms behind intelligence at these levels. They
provide strong theoretical justification and computational evidence
for the claim:
\begin{quote}
    \centering
    \textit{Machines can learn empirical knowledge and conduct
    inductive inference.}
\end{quote}
So far, however, no rigorous theoretical or scientific evidence
suggests that these mechanisms alone would suffice to achieve the
high-level human intelligence that the original AI program truly
aimed to understand and wanted for machines to imitate or even surpass.

In fact, as of today, we know little about how to rigorously verify
or certify whether a system is truly capable of such high-level
intelligence, even though the Turing Test, also known as the
Imitation Game,\footnote{In Turing's proposal, the evaluator, or
    interrogator, is a human. However, most human evaluators have limited
    scientific training and knowledge, and their conclusions can be
subjective.} was proposed in 1950 \citep{Turing-1950}, as an initial
attempt to address the fundamental question:
\begin{quote}
    \centering
    \textit{Can machines think?}
\end{quote}
For a long time, a precise definition of such a test was not deemed
necessary since machine capabilities were far below those of humans
(or even animals). However, given recent technological advances, many
models and systems now claim to reach or even surpass human
intelligence. Therefore, it is high time to develop a scientific and
executable definition of the Turing test, i.e., a systematic and
rigorous protocol to evaluate and certify the intelligence level of a
model or a system that claims to be ``intelligent.''

As Turing argued in his original paper, it is rather difficult
to define what ``thinking'' is. We may instead start with something more
tangible and try to determine whether a system can truly
``understand'' certain concepts or knowledge. For example, how can we
rigorously verify whether an intelligent model or system has truly
understood an abstract concept such as:
\begin{enumerate}
    \item the notion of equality between two quantities,
    \item the notion of numbers (natural, rational, real, or imaginary),
    \item the notion of infinity and mathematical induction,
    \item the notion of logic consistency\footnote{In mathematics, if a set of axioms and their deduced results are consistent and 
    do not contradict one another, then they exist. This is known as Hilber's Principle. This is different from the 
    notion of ``consistency'' in Chapter \ref{ch:consistent-representations}, which requires the learned
    representations to be consistent with the empirical data distributions. One notion of consistency is inductive in nature 
    and the other is deductive.} and proof by contradition,
\end{enumerate}
just to name a few? Or has it simply memorized a massive number of
examples of such notions? Note that state-of-the-art large language
models still struggle with simple mathematical questions like: ``Is
3.11 larger or smaller than 3.9?''\footnote{Some models have
    corrected their answers to such questions through targeted
    engineering, or have incorporated additional reasoning mechanisms
    that verify immediate answers and correct them during inference.
    However, we leave it to the reader as an exercise to rigorously test
    whether any state-of-the-art language model truly understands the
    notion of numbers (natural, rational, real, and complex) and their
associated arithmetic.}

From this book, we know modern generative AI systems primarily rely
on Bayesian (hence inductive) inference to produce their answers. Can
machines also truly grasp mathematical induction or general logical
deduction and understand its difference from Bayesian induction?
How do we verify whether a system truly understands the rules of
logical and mathematical deduction and can apply them rigorously? Or,
again, has it merely memorized a large number of instances of such
deductions, such as chain-of-thought data, and learned and
exploited their distributional patterns? Hence we have an open question:
\begin{quote}
    \centering
    \textit{Is there any difference between \textbf{memorizing} and
    \textbf{understanding}?}
\end{quote}
Only by knowing the difference can we possibly attempt to provide a
meaningful answer to the question: ``\textit{Can machines understand?}''

Probably much more importantly, how do we verify whether a self-claimed
intelligent system is capable of \textit{creating abstract concepts or improving its own knowledge} proactively via deductive means? For example, can it create new mathematical
concepts and discover new physical laws or causal relationships?
The connection between understanding and creating can be rather fundamental, as the physicist Richard Feynman famously put it:
\begin{quote}
    \centering
    \textit{``What I cannot create, I do not understand.''}
\end{quote}


\subsection{Scientific Tests of Intelligence}
Hence, it is high time we develop scientifically sound evaluation
methods to determine which of the following categories a system's
intelligent capability belongs to:
\begin{enumerate}[label={Level-\arabic*.}, leftmargin=*]
    \item having merely learned the distribution of some
        knowledge-carrying data and being able to regenerate them;
    \item being capable of autonomously and continuously
        learning new knowledge or correcting existing knowledge from new
        observations and experiences;
    \item being able to truly understand and create abstract
        concepts, rules and laws and knowing how to apply them
        correctly in a deductive manner;
    \item being able to generate new scientific hypotheses or math conjectures and verify or falsify them based on experimentation or logical deduction.
\end{enumerate}

\begin{figure}[t]
    \centering
    \includegraphics[width=\linewidth]{\toplevelprefix/chapters/chapter9/figs/tests.png}
    \caption{Three tests for different levels or types of
        intelligence capabilities: the Wiener test for basic
        intelligence, the Turing test for human-level intelligence, and
    the Popper test for scientist-level intelligence.}
    \label{fig:three-tests}
\end{figure}

To evaluate and distinguish these different levels of intelligence,
we suggest that there should probably be at least three different
tests, as illustrated in \Cref{fig:three-tests}:
\begin{enumerate}
    \item \textit{The Norbert Wiener Test:} to determine whether a
        system can self-correct and create new (empirical) knowledge
        on its own or merely receives and updates information
        passively through external reinforcement or supervision;
    \item \textit{The Alan Turing Test:} to determine whether a
        system can create and understand abstract concepts and
        knowledge or merely learns and memorizes statistics of the
        samples and uses them for Bayesian inference;
    \item \textit{The Karl Popper Test:} to determine whether a
        system can proactively explore and create new abstract
        knowledge by forming and verifying hypotheses or theories
        based on experimental consistency or logical self-consistency.
\end{enumerate}
We believe that, for such tests, the evaluator, or interrogator,
should not be a single human or an arbitrary group of humans but rather a
scientifically certified entity following a scientifically sound
protocol and process so that the conclusions would be rigorous and trustworthy.

As we have seen throughout this book, \textit{compression} has played
a fundamental role in developing a memory or creating empirical knowledge. It is the governing principle and a universal mechanism for identifying an
empirical data distribution and organizing the information encoded therein with a structured representation. To a large extent, it explains the practice of ``artificial intelligence'' with deep networks, which is largely inductive in nature.\footnote{Systems that have built-in deductive mechanisms for conducting formal verification are not in this category.} An outstanding question for future
study is whether \textit{compression alone} is sufficient to achieve
all the higher-level deductive intelligence mentioned above.
In particular, are mathematical abstraction, causal inference, hypothesis
generation, and logical reasoning and deduction some kind of extended or
transcended forms of compression? Is there some fundamental difference between learning data distributions through compression and forming high-level concepts and theories through abstraction? Hence we have another open question:
\begin{quote}
    \centering
    \textit{Is there any difference between \textbf{compression} and
    \textbf{abstraction}?}
\end{quote}
Only by knowing the difference can we possibly attempt to provide a meaningful answer to the question: ``\textit{Can machines create knowledge proactively via deduction?}''

To a large extent, Science---and its associated code book,
Mathematics---can be viewed as the most advanced form of
intelligence, unique to educated and enlightened humans. Philosopher Sir Karl Popper once suggested:
\begin{quote}
    \centering
    \textit{``Science may be described as the art of systematic
    oversimplification.''}
\end{quote}
We believe that uncovering and understanding the underlying
mathematical principles and computational mechanisms of such higher-level intelligence, and successfully reproduce them through machines,
 will be the final frontier for Science, Mathematics, and Computation altogether!

\end{document}
